\chapter{\texttt{Entity}, \texttt{Relation} y los tipos públicos}
\label{cap:entity-relation}

\begin{intuicion}
El modelo de datos público de GraphHunter ---el que ven los parsers,
los \emph{transports} y la API--- es deliberadamente legible: un
\texttt{Entity} es un nodo con identificador humano y tipo, un
\texttt{Relation} es una arista con origen, destino, tipo y
\emph{timestamp}. Los dos llevan un \texttt{HashMap} de
\texttt{metadata} opcional. Esa forma es cómoda para el ingestor y
los parsers, pero \emph{no} es la forma en la que el matcher itera
en su \emph{hot loop}; ahí entra \texttt{CompactRelation}
(capítulo~\ref{cap:compact-relation}). Este capítulo presenta el
modelo público desde el evento de Sysmon hasta los triples que el
ingestor empuja al grafo.
\end{intuicion}

\section{Tipos públicos}

\begin{lstlisting}[language=Rust]
pub struct Entity {
    pub id: String,                              // ej. "10.0.0.42" o "lsass.exe"
    pub entity_type: EntityType,                 // ej. IP, Process, User
    pub metadata: HashMap<String, MetadataValue>, // campos opcionales del log
}

pub struct Relation {
    pub source_id: String,
    pub dest_id: String,
    pub rel_type: RelationType,                  // ej. Auth, Connect, Spawn
    pub timestamp: i64,                           // segundos epoch
    pub metadata: HashMap<String, MetadataValue>,
}
\end{lstlisting}

Tres observaciones:

\begin{enumerate}
    \item \texttt{id} es un \texttt{String} humano (\texttt{"alice"},
    \texttt{"db-01"}, \texttt{"10.0.0.42"}). Es la \emph{primary
    key} del vértice; dos eventos con la misma \texttt{id} y
    \texttt{entity\_type} se refieren al mismo nodo.
    \item \texttt{entity\_type} es un \texttt{enum}: nueve variantes
    fijas (\texttt{User}, \texttt{Host}, \texttt{Process},
    \texttt{IP}, \texttt{File}, \texttt{Domain}, \texttt{Service},
    \texttt{Resource}, \texttt{Group}). No es extensible en runtime;
    cada nuevo tipo requiere agregar la variante y recompilar.
    \item \texttt{metadata} es flexible. El parser de Sysmon mete acá
    \texttt{ProcessGuid}, \texttt{CommandLine},
    \texttt{ImagePath}; el parser de Sentinel mete
    \texttt{TenantId}, \texttt{ResourceGroup}. El matcher lo consulta
    solamente cuando hay un predicado en la hipótesis (capítulo
    \ref{cap:dsl-chain}).
\end{enumerate}

\section{\texttt{EntityType} y \texttt{RelationType}: alfabetos
cerrados}

\begin{lstlisting}[language=Rust]
pub enum EntityType {
    User      = 0,
    Host      = 1,
    Process   = 2,
    IP        = 3,
    File      = 4,
    Domain    = 5,
    Service   = 6,
    Resource  = 7,
    Group     = 8,
    Other     = 255,
}
\end{lstlisting}

Los \emph{tags} numéricos (0--8) son explícitos por dos razones:

\begin{itemize}
    \item Permiten almacenar el tipo en un \texttt{u8} (1 byte) en
    vez de un \texttt{String} de 6--8 bytes.
    \item Permiten indexar arrays de tamaño fijo
    \texttt{[u32; 9]} en NLF (capítulo~\ref{cap:nlf-detalle}) sin
    \emph{branch} en el lookup.
\end{itemize}

\(\texttt{Other = 255}\) es un \emph{escape hatch} para casos
ambiguos del parser; el motor lo cuenta separadamente y NO lo
incluye en NLF (sería ruido).

\texttt{RelationType} es similar, con $\sim$15 variantes
(\texttt{Auth}, \texttt{Connect}, \texttt{Spawn}, \texttt{Read},
\texttt{Write}, \texttt{Beacon}, \texttt{Resolve}, \texttt{Member},
etc.). El motor usa el \texttt{u8 tag} en
\texttt{CompactRelation::rel\_type\_tag} para matching sin resolver
el enum entero.

\section{De evento Sysmon a triples}

Veamos cómo un evento concreto se convierte en \texttt{Entity}s y
\texttt{Relation}s.

\begin{ejemplo}[Sysmon Event ID 1: Process Create]
Un evento Sysmon NDJSON crudo:

\begin{lstlisting}
{
    "EventID": 1,
    "TimeCreated": "2026-04-15T14:23:51Z",
    "Computer": "WS-001",
    "User": "CORP\\\\m.gomez",
    "Image": "C:\\\\Windows\\\\System32\\\\powershell.exe",
    "ParentImage": "C:\\\\Windows\\\\explorer.exe",
    "ProcessGuid": "{ABCD-1234-...}",
    "ParentProcessGuid": "{EFGH-5678-...}",
    "CommandLine": "powershell.exe -enc ZQB4ZQ==",
    "Hashes": "MD5=abc...,SHA256=def..."
}
\end{lstlisting}

El parser \filepath{platform/parsers/sysmon} produce de este evento:

\textbf{3 entidades} (idempotente: si ya existen, se ignoran):
\begin{itemize}
    \item \texttt{Entity\{ id="CORP\textbackslash{}m.gomez", entity\_type=User, metadata=\{\}\}}
    \item \texttt{Entity\{ id="WS-001", entity\_type=Host, metadata=\{\}\}}
    \item \texttt{Entity\{ id="\{ABCD-1234-...\}", entity\_type=Process,\\
              metadata=\{"image": "powershell.exe", "cmdline": "...", "md5": "abc...", "sha256": "def..."\}\}}
    \item \texttt{Entity\{ id="\{EFGH-5678-...\}", entity\_type=Process,\\
              metadata=\{"image": "explorer.exe"\}\}}
\end{itemize}

\textbf{2 relaciones}:
\begin{itemize}
    \item \texttt{Relation\{ source="WS-001", dest="\{ABCD-1234-...\}", rel\_type=Hosts, ts=1744732431, metadata=\{\}\}}
    \item \texttt{Relation\{ source="\{EFGH-5678-...\}", dest="\{ABCD-1234-...\}", rel\_type=Spawn, ts=1744732431, metadata=\{\}\}}
\end{itemize}

Notar:
\begin{itemize}
    \item Un solo evento generó \emph{dos triples} (proceso hijo
    spawn-eado, host que lo aloja). Los parsers son ``\emph{event-to-N-triples}'',
    no ``\emph{one-to-one}''.
    \item Los IDs de los procesos son los \texttt{ProcessGuid}
    (estables), no los \texttt{ProcessId} (PIDs reciclados por
    Windows).
    \item El \texttt{TimeCreated} ISO-8601 se convierte a
    \texttt{i64} segundos epoch ---el motor jamás guarda timestamps
    como string.
\end{itemize}
\end{ejemplo}

\section{El contrato del parser}

Cualquier \emph{adapter} de fuente nueva tiene que cumplir:

\begin{lstlisting}[language=Rust]
pub trait LogParser {
    type RawRow;

    fn parse(&self, row: Self::RawRow)
        -> Result<Vec<ParsedTriple>, ParseError>;
}

pub struct ParsedTriple {
    pub source_id: String,
    pub source_type: EntityType,
    pub dest_id: String,
    pub dest_type: EntityType,
    pub rel_type: RelationType,
    pub timestamp: i64,
    pub metadata: HashMap<String, MetadataValue>,
}
\end{lstlisting}

\noindent
Un \texttt{ParsedTriple} es ``\emph{Entity} + \emph{Entity} +
\emph{Relation}'' aplanados en una sola estructura ---los IDs y los
tipos identifican unívocamente las dos entidades involucradas, así
que no hace falta enviarlas por separado. El \emph{ingestor}
desempaqueta cada triple en dos \texttt{add\_entity} (idempotente)
más un \texttt{add\_relation}.

\begin{enelcodigo}
\filepath{platform/parsers/src/lib.rs} define el trait. Cada parser
concreto vive en su sub-módulo: \texttt{sysmon}, \texttt{sentinel},
\texttt{evtx}, \texttt{cognito}, \texttt{forti}, \texttt{iis},
\texttt{generic\_csv}.
\end{enelcodigo}

\section{Idempotencia y deduplicación}

Una propiedad importante del modelo: ingerir el mismo log dos veces
no debería duplicar entidades. Si ``alice'' ya existe como
\texttt{User}, el segundo \texttt{add\_entity("alice", User, ...)}
es no-op (\emph{merge} de metadata).

Esto permite reconstrucciones idempotentes: si el ingestor crashea a
mitad de un dataset, reanudarlo desde el principio no genera datos
duplicados. La idempotencia también vale para \texttt{add\_relation},
con una salvedad: dos eventos con la \emph{misma} terna
\((u, v, \text{rel\_type})\) en \emph{distinto} timestamp son
relaciones distintas (es un multigrafo). El motor las almacena ambas.

\section{Metadata: cuándo materializa, cuándo no}

El \texttt{HashMap<String, MetadataValue>} de cada entidad y
relación es flexible pero costoso. El motor aplica una optimización:
si nadie consulta la metadata, no se materializa en RAM.

\begin{lstlisting}[language=Rust]
pub enum MetadataValue {
    Str(String),
    Int(i64),
    Float(f64),
    Bool(bool),
    Arr(Vec<MetadataValue>),
    Obj(HashMap<String, MetadataValue>),
}
\end{lstlisting}

En el \texttt{CompactRelation} (capítulo~\ref{cap:compact-relation}),
la metadata se almacena por \emph{offset} a un
\texttt{MetadataStore} append-only; un offset de 0 significa ``sin
metadata'' y se chequea sin lookup. Los hunts que no tienen
predicados sobre metadata jamás resuelven el offset; el costo es
amortizado.

\section{Resumen}

\begin{itemize}
    \item Modelo público: \texttt{Entity} (id + tipo + metadata) y
    \texttt{Relation} (origen + destino + tipo + timestamp +
    metadata). Legible, parser-friendly.
    \item Tipos cerrados (9 \texttt{EntityType}, $\sim$15
    \texttt{RelationType}) con tag numérico explícito que
    \texttt{CompactRelation} reusa.
    \item Parsers son \emph{event-to-N-triples}; el ingestor
    desempaqueta cada triple en \texttt{add\_entity} (idempotente)
    más \texttt{add\_relation} (append).
    \item Metadata se materializa lazy via \texttt{MetadataStore}
    append-only; offset 0 significa ``sin metadata''.
\end{itemize}

\begin{ejercicios}
\begin{enumerate}
    \item \starred{} Para el evento Sysmon Event ID 3 (Network
    Connection): listar qué entidades y qué relaciones genera el
    parser. (Pista: hay un proceso, un host, una IP destino, un
    puerto.)

    \item ¿Por qué \texttt{EntityType} tiene tag numérico explícito
    (\texttt{User = 0}, \texttt{Host = 1}, etc.)? Imaginá la
    alternativa de dejar que Rust elija; ¿qué problemas surgen
    para FFI y serialización?

    \item Imaginá que mañana aparece una nueva fuente cuyo modelo
    natural NO encaja en los 9 \texttt{EntityType} actuales (digamos,
    ``contenedores Kubernetes''). Discutí dos diseños: (a) agregar
    \texttt{EntityType::Container = 9} y recompilar; (b) usar
    \texttt{Other} y poner el tipo real en metadata. Trade-offs.

    \item \starred{} El \texttt{HashMap<String, MetadataValue>} usa
    strings como keys, lo que cuesta memoria. Diseñá una alternativa
    con un esquema fijo por \texttt{EntityType} (struct estático).
    ¿Qué ganás? ¿Qué perdés?

    \item Discutí: si ingerís el mismo evento dos veces, ¿se duplica
    la \texttt{Relation}? ¿Por qué la respuesta correcta NO es
    siempre ``no''? (Pista: dos eventos con el mismo timestamp pero
    distinta metadata son distintos eventos del log, aunque
    \emph{conceptualmente} se refieran a lo mismo.)
\end{enumerate}
\end{ejercicios}
